\documentclass[conference]{IEEEtran}

\usepackage{amsmath}
\usepackage{booktabs}
\usepackage{graphicx}
\usepackage{microtype}
\usepackage{url}
\usepackage{xcolor}
\usepackage{array}
\usepackage[hidelinks]{hyperref}

\newcommand{\system}{AssetMind AI}
\newcommand{\todo}[1]{\textcolor{red}{[REPLACE: #1]}}

\title{\system: A Tool-Calling, Evidence-Traceable Agent for Industrial Asset Fault Investigation}

\author{
\IEEEauthorblockN{\todo{Student Name}}
\IEEEauthorblockA{
\todo{Department or Programme}\\
IILM University\\
\todo{City, Country}\\
\todo{institutional email}}
}

\begin{document}
\maketitle

\begin{abstract}
Industrial fault investigation requires engineers to combine heterogeneous
evidence, including sensor observations, historical faults, maintenance
records, and operating procedures. Conventional retrieval-augmented
generation (RAG) systems frequently implement retrieval as a fixed pipeline
and provide limited assurance that a generated claim is supported by the
evidence used for that interaction. This paper presents \system, a
full-stack generative-artificial-intelligence fault-investigation system in
which a language model selects among six typed tools: semantic document
retrieval, keyword document retrieval, asset context, fault history,
maintenance history, and a deterministic guardrail engine. The system is
designed as an agent rather than a fixed retrieve-then-generate pipeline.
Every tool call and result is persisted as versioned evidence, and answers use
machine-checkable citations verified against evidence from the same turn.

The implementation combines Next.js, the Vercel AI SDK, PostgreSQL 16 with
pgvector and PostgreSQL full-text search, Prisma, OpenAI models, Clerk
authentication, Amazon Web Services (AWS) RDS, AWS S3, and Inngest. The
guardrail engine implements deterministic HDF, PWF, and OSF rules of the AI4I
2020 dataset. These rules reproduce their labelled rows exactly: 115 HDF
rows, 95 PWF rows, and 98 OSF rows. Tool-wear failure and random failure are
not represented as deterministic thresholds because the source data do not
permit that inference.

Implementation validation passed: 10,000 dataset rows were checked, 66
document chunks were embedded with \texttt{text-embedding-3-small}, both
keyword and vector retrieval returned provenance-bearing results, 202
automated tests passed, and the production health endpoint reported a live
database connection. A frozen evaluation protocol defines 24 labelled
questions, five tool-subset configurations, and three repeats (360 runs).
Full comparative agent results are intentionally not claimed until that grid
is executed. The contribution is therefore a reproducible agent,
guardrail, evidence, and evaluation design with verified implementation
correctness, rather than an unsupported claim of comparative model
superiority.
\end{abstract}

\begin{IEEEkeywords}
industrial fault investigation, tool-calling agents, retrieval-augmented
generation, predictive maintenance, pgvector, traceability, AI4I 2020
\end{IEEEkeywords}

\section{Introduction}
Industrial equipment investigation is not a single retrieval problem. An
engineer may need to identify an asset, inspect a sensor reading, compare it
with previous failures, consult a maintenance procedure, and determine
whether a rule-based signal is relevant. These sources differ in structure,
units, temporal meaning, and evidential authority. A fluent answer without a
verifiable evidence trail is consequently difficult to trust and evaluate.

Large language models (LLMs) make conversational investigation interfaces
possible, but introduce three risks. First, the model may make unsupported
claims when a requested fact is absent. Second, retrieval may be implemented
as a hidden fixed pipeline, making it unclear whether an answer benefited
from semantic search, exact search, structured records, or a rule engine.
Third, a citation displayed in an answer may not correspond to evidence
returned during that interaction.

\system addresses these risks through four design decisions:
\begin{enumerate}
\item The model receives a subset of a common tool registry and chooses
which tools to call and in what order; no retrieve-then-generate sequence is
forced.
\item Tool calls, outcomes, and citable results are persisted in a versioned
evidence structure attached to the chat message.
\item Answer citations are checked against evidence produced by the same
turn rather than trusted by visual inspection.
\item The AI4I rules are implemented as structured, auditable verdicts
containing observed values, thresholds, and per-clause explanations. Modes
that cannot be inferred deterministically are not assigned invented
thresholds.
\end{enumerate}

The study defines five research questions (RQs): RQ1 asks how document
retrieval changes groundedness over a no-tool baseline; RQ2 asks what
structured asset, fault, and maintenance tools add; RQ3 asks whether the
guardrail tool is used appropriately; RQ4 asks about cost and latency; and
RQ5 asks whether factual claims can be traced mechanically to same-turn
evidence. The contributions are: 1) a deployable tool-calling
fault-investigation architecture, 2) a versioned evidence and citation
verification design, and 3) a reproducible evaluation protocol that separates
verified implementation results from pending comparative experiments.

\section{Related Work}
\subsection{Retrieval-Augmented Generation}
RAG combines a generator with an external non-parametric memory
\cite{lewis2020rag}. It can reduce reliance on knowledge encoded in model
parameters and make answers responsive to a changing document collection.
However, retrieval alone does not establish that the model used a result,
that the result was relevant, or that each factual sentence is supported.
\system therefore evaluates both retrieval availability and the evidence
trace emitted by the agent.

\subsection{Reasoning and Acting with Tools}
ReAct demonstrated the value of interleaving reasoning and actions through
tools \cite{yao2023react}. \system adopts the tool-use principle without
introducing a separate agent framework. The Vercel AI SDK executes the
multi-step loop, while a typed registry determines which tools are
available. This permits ablation without duplicating implementation paths.

\subsection{Industrial Predictive Maintenance}
The AI4I 2020 dataset provides 10,000 simulated industrial-process rows with
temperature, speed, torque, tool-wear, product-quality, machine-failure, and
failure-mode fields \cite{matzka2020ai4i}. It is useful because it is public,
labelled, and reproducible. It is not a pump dataset. The pump-style asset
identifiers in this project are a declared plant-role abstraction; the
underlying readings and rules retain their milling-machine meaning. The
system does not claim to model pump physics.

\subsection{Traceability}
Displaying a source label is weaker than proving that the labelled source
was returned during the interaction. \system treats traceability as an
observable property: a citation must resolve to an evidence item recorded
for that turn. This catches invented identifiers, wrong chunk ordinals, and
guardrail citations for modes not returned by the tool. Mechanical validity
does not prove semantic entailment, so the protocol supplements it with
sentence coverage and human review.

\section{Research Questions and Experimental Design}
The evaluation is a controlled tool-availability ablation. The model,
temperature, system-prompt version, output limits, question set, pricing,
and repeat count are frozen before the agent experiments:
\begin{itemize}
\item \textbf{RQ1:} How much does document retrieval improve groundedness
over a no-tool baseline?
\item \textbf{RQ2:} What value do structured asset, fault, and maintenance
tools add beyond document retrieval?
\item \textbf{RQ3:} Does the agent call guardrails for relevant readings and
avoid unnecessary guardrail claims?
\item \textbf{RQ4:} What cost and latency trade-off accompanies additional
tool availability?
\item \textbf{RQ5:} Can factual claims be traced to same-turn tool evidence?
\end{itemize}

The five frozen configurations are shown in Table~\ref{tab:configs}.

\begin{table}[t]
\caption{Frozen tool-subset configurations}
\label{tab:configs}
\centering
\footnotesize
\begin{tabular}{c p{0.67\columnwidth}}
\toprule
Config. & Registered tools\\
\midrule
1 & None (baseline)\\
2 & Vector document retrieval\\
3 & Vector and keyword document retrieval\\
4 & Configuration 3 plus asset context, fault history, and maintenance history\\
5 & Configuration 4 plus deterministic guardrails\\
\bottomrule
\end{tabular}
\end{table}

The labelled set contains 24 questions: eight retrieval, four structured,
six guardrail, three mixed, two weak-signal, and one abstention question.
Each question is evaluated in every configuration and repeated three times,
for $24\times5\times3=360$ runs. The model is \texttt{gpt-4o-mini} at
temperature zero. Runs record the system-prompt version and Git commit.

\section{System Architecture}
\subsection{Application and Data Topology}
\system is a single Next.js application. Authenticated users access a
dashboard, asset pages, monitoring views, and an investigation chat. Server
Components read dashboard data directly through Prisma. The chat route is a
streaming endpoint because it executes the LLM loop and returns incremental
UI messages. The deployed topology is:
\begin{center}
\footnotesize
\texttt{Browser $\rightarrow$ Vercel/Next.js $\rightarrow$
OpenAI + AWS RDS + S3/Inngest}
\end{center}

PostgreSQL 16 stores structured assets, components, sensor readings, fault
records, maintenance logs, documents, chunks, chats, messages, and daily
usage. pgvector stores 1536-dimensional embeddings; a generated
\texttt{tsvector} column and GIN index support keyword search. The deployed
database contains eight assets, 64 components, 10,000 readings, 330
attributable fault records, 104 maintenance logs, 32 documents, and 66
chunks.

\subsection{Tool Registry}
The six tools are:
\begin{itemize}
\item \texttt{searchDocumentsVector}: semantic document search;
\item \texttt{searchDocumentsKeyword}: exact PostgreSQL full-text search;
\item \texttt{getAssetContext}: identity, components, and current state;
\item \texttt{getFaultHistory}: historical faults and outcomes;
\item \texttt{getMaintenanceHistory}: maintenance and inspection records;
\item \texttt{checkGuardrails}: deterministic HDF, PWF, and OSF verdicts.
\end{itemize}

\texttt{buildToolRegistry(names)} is the only registry-construction entry
point used by the user interface, chat route, and evaluation harness. The
five configurations are subsets of this registry, not five pipelines. Tool
inputs use Zod schemas, and tool failures return typed outcomes rather than
terminating the stream. Step, token, timeout, retry, quota, result-cap, and
evidence-budget limits bound the runtime.

\subsection{Prompt, Evidence, and Citations}
The system prompt is versioned as \texttt{v1} and held constant across
configurations. It requires every factual claim to be supported by same-turn
tool evidence, instructs the model to abstain when evidence is unavailable,
and treats document content as untrusted data rather than instructions.
Citations use the following machine-checkable forms:
\begin{itemize}
\item \texttt{[chunk:<id>\#<ordinal>]} for document chunks;
\item \texttt{[fault:<id>]} for fault records;
\item \texttt{[guardrail:<mode>]} for an observed HDF, PWF, or OSF verdict.
\end{itemize}

Evidence is persisted in \texttt{ChatMessage.evidence} with a version field.
The verifier extracts citations, resolves each identifier against same-turn
evidence, and reports valid and invalid citations. This provides a mechanical
traceability check while acknowledging that semantic entailment still
requires review.

\section{Deterministic Guardrail Engine}
The guardrail engine is a rules component, not a trained classifier. For a
sample with air temperature $T_a$, process temperature $T_p$, speed $s$,
torque $\tau$, tool wear $w$, and quality type $q$, the implemented rules
are:
\begin{align}
\mathrm{HDF} &\iff |T_p-T_a|<8.6\ \land\ s<1380,\\
P &= \tau\frac{2\pi s}{60},\\
\mathrm{PWF} &\iff P<3500\ \lor\ P>9000,\\
\mathrm{OSF} &\iff w\tau>L_q,
\end{align}
where $L_L=11000$, $L_M=12000$, and $L_H=13000$.
Temperatures are stored in Kelvin, torque in Nm, speed in rpm, and tool wear
in minutes.

The formulas were verified row-by-row against the source CSV before being
used by the engine. As shown in Table~\ref{tab:rules}, all three
deterministic modes have exact row-set agreement.

\begin{table}[t]
\caption{Guardrail verification against AI4I labels}
\label{tab:rules}
\centering
\footnotesize
\begin{tabular}{lrrl}
\toprule
Mode & Label rows & Computed rows & Agreement\\
\midrule
HDF & 115 & 115 & Exact\\
PWF & 95 & 95 & Exact\\
OSF & 98 & 98 & Exact\\
\bottomrule
\end{tabular}
\end{table}

TWF and RNF are not deterministic functions of the stored process variables.
For TWF, 790 rows fall in the 200--240 minute wear band while only 46 are
labelled TWF, and three TWF rows lie outside the band. RNF is a per-row random
event; the released CSV contains 19 such labels. The system therefore reports
these modes as irreducibly random instead of assigning scientifically
unsupported thresholds.

\section{Materials and Methods}
\subsection{Dataset and Corpus}
AI4I 2020 is used as the sensor and ground-truth source
\cite{matzka2020ai4i}. The seeded application maps rows to eight named assets
using a fixed seed. Synthetic manuals, SOPs, fault narratives, and
inspection reports provide a reproducible 32-document, 66-chunk corpus.
The documents are explicitly synthetic and the asset naming abstraction is
explicitly acknowledged as a limitation.

\subsection{Embedding and Retrieval}
Chunks are embedded with \texttt{text-embedding-3-small} at 1536 dimensions.
Vectors are normalized and searched with pgvector cosine distance. Keyword
retrieval searches the generated \texttt{tsvector} column using a GIN index.
The evaluation uses exact vector scans because the corpus is small; this
avoids approximate-nearest-neighbour parameters becoming an uncontrolled
variable. Both retrieval paths return source identifiers and ordinals.

\subsection{Evaluation Metrics}
For RQ1, groundedness is sentence citation coverage: the fraction of answer
sentences containing a citation resolving to same-turn evidence. RQ2
compares chunk, fault, and guardrail evidence composition and tool-call mix.
RQ3 measures relevant guardrail calls, deterministic mode agreement, and
irrelevant over-calling; TWF and RNF are excluded from threshold
precision/recall. RQ4 records all-step usage, final-step usage, calculated
cost, and wall-clock duration. RQ5 reports citation validity, sentence
coverage, and human review of uncited factual claims.

For citations:
\begin{equation}
\mathrm{validity} =
\frac{\mathrm{valid\ citations}}
{\mathrm{valid\ citations}+\mathrm{invalid\ citations}}.
\end{equation}
The no-citation case is not treated as an invalid citation; it is measured by
sentence coverage.

\subsection{Reproducibility}
The protocol and question set are version-controlled and frozen. The
repository provides database, guardrail, retrieval, citation, and evaluation
scripts. Continuous integration runs linting, typechecking, the guardrail
dataset check, unit tests, and a production build. Production uses Vercel,
AWS RDS, AWS S3, OpenAI, Clerk, and Inngest. The deployed RDS instance has
one-day backups and CloudWatch alarms for CPU, storage, and connections.

\section{Results}
\subsection{Implementation Validation}
Table~\ref{tab:validation} summarizes results directly verified by tests,
read-only database checks, or deployment checks.

\begin{table}[t]
\caption{Verified implementation results}
\label{tab:validation}
\centering
\footnotesize
\begin{tabular}{p{0.62\columnwidth}r}
\toprule
Validation & Result\\
\midrule
Automated test files & 20 passed\\
Automated tests & 202 passed\\
Dataset rows checked & 10,000\\
Embedded chunks & 66/66\\
Keyword results & 5 with provenance\\
Vector results & 5 with provenance\\
Question-set questions & 24 valid\\
Production health & HTTP 200; database up\\
CI for commit \texttt{c2b9743} & Passed\\
\bottomrule
\end{tabular}
\end{table}

The live retrieval check returned five keyword and five vector results. Vector
distances were finite, within the cosine-distance range, and sorted in
ascending order. A nonzero nearest-neighbour distance is expected for a
free-text query that does not exactly equal a stored chunk.

\subsection{Scope of the Current Results}
The frozen protocol and sample manifests exist, but the complete 360-run
agent grid and planned human review are not included in the verified results
reported here. Accordingly, this paper does not claim a population-level
groundedness gain, structured-tool improvement, guardrail-use rate,
cost/latency curve, or complete unsupported-claim rate. Those findings must
be added only after the full grid is executed and its manifests and human
adjudication are preserved.

\section{Discussion}
The architecture is appropriate for investigation because questions require
different evidence paths. A procedure question may require document search,
a historical question may require relational tools, a concrete reading may
require guardrails, and an absent fact should produce abstention. A common
registry makes these alternatives measurable without code duplication.

The guardrail analysis also demonstrates why labelled fields must not
automatically be treated as predictable targets. HDF, PWF, and OSF are
reproducible under verified compound rules; TWF and RNF are not. Treating all
five modes as thresholds would create a seemingly complete but scientifically
incorrect system.

\section{Limitations and Threats to Validity}
\textbf{Construct validity:} sentence citation coverage is a proxy for
groundedness. A cited sentence can still overstate its source, while a
correct sentence can be uncited; human review is required.

\textbf{Internal validity:} tool availability changes the context supplied to
the model. Configuration 5 may provide more evidence and tokens than
configuration 1, so a gain cannot be attributed solely to one retrieval
method. The protocol declares this context-budget confound.

\textbf{External validity:} the sensor data come from one simulated
milling-machine dataset, the documents are synthetic, and pump-style names
are an abstraction. The results must not be generalized to real pump physics
or safety-critical deployment.

\textbf{Conclusion validity:} three repeats and a 24-question set are
appropriate for a capstone prototype but limit statistical power. Temperature
zero does not remove provider-side or infrastructure variance.

\textbf{Operational security:} the current RDS security group permits public
PostgreSQL access on port 5432. A production deployment should use private
networking, a verified fixed-egress solution, or a narrowly scoped allowlist.
The limitation is documented rather than hidden.

\section{Ethical and Safety Considerations}
\system is a decision-support system, not an autonomous controller. It does
not issue machine commands or replace an engineer's inspection. It abstains
when evidence is unavailable and distinguishes deterministic guardrails from
irreducible random modes. Clerk scopes chat records by user, runtime limits
bound tool and token use, and document text is treated as untrusted data to
reduce prompt-injection risk. Dataset provenance and license conditions must
be retained in the final submission.

\section{Conclusion}
\system demonstrates a reproducible architecture for evidence-traceable
industrial fault investigation. Its central design is a tool-calling agent
whose available actions can be ablated without creating separate pipelines.
Structured records, two retrieval paths, and deterministic guardrails are
exposed through one registry. Evidence is captured during execution and
citations are verified against the same-turn trace.

The implementation is validated at software and data levels: the
deterministic guardrails reproduce HDF, PWF, and OSF labels exactly;
database, pgvector, full-text search, and embeddings are operational;
automated tests and CI pass; and the deployed health endpoint reaches the
live database. Comparative agent performance is intentionally left as a
measured future result until the frozen 360-run evaluation is complete.

\section*{Acknowledgment}
The authors thank \todo{supervisor or laboratory, if applicable} for
guidance and review.

\begin{thebibliography}{00}
\bibitem{matzka2020ai4i}
S.~Matzka, ``Explainable artificial intelligence for predictive maintenance
applications,'' in \emph{Proc. 3rd Int. Conf. Artificial Intelligence for
Industries (AI4I)}, 2020, pp.~1--8, doi:
\href{https://doi.org/10.1109/AI4I49448.2020.00023}{10.1109/AI4I49448.2020.00023}.

\bibitem{lewis2020rag}
P.~Lewis, E.~Perez, A.~Piktus, F.~Petroni, V.~Karpukhin, N.~Goyal,
H.~Küttler, M.~Lewis, W.-t.~Yih, T.~Rocktäschel, S.~Riedel, and
S.~Kiela, ``Retrieval-augmented generation for knowledge-intensive NLP
tasks,'' in \emph{Advances in Neural Information Processing Systems},
vol.~33, 2020, pp.~9459--9474.

\bibitem{yao2023react}
S.~Yao, J.~Zhao, D.~Yu, N.~Du, I.~Shafran, K.~Narasimhan, and Y.~Cao,
``ReAct: Synergizing reasoning and acting in language models,'' in
\emph{Proc. Int. Conf. Learning Representations}, 2023.

\bibitem{johnson2019faiss}
J.~Johnson, M.~Douze, and H.~Jégou, ``Billion-scale similarity search with
GPUs,'' \emph{IEEE Trans. Big Data}, vol.~7, no.~3, pp.~535--547, 2021,
doi: 10.1109/TBDATA.2019.2921572.

\bibitem{postgresqlfts}
PostgreSQL Global Development Group, ``Full text search,'' PostgreSQL
Documentation, 2026. [Online]. Available:
\url{https://www.postgresql.org/docs/current/textsearch.html}

\bibitem{pgvector}
P.~Garrison and A.~Ramirez, ``pgvector: Open-source vector similarity search
for Postgres,'' 2026. [Online]. Available:
\url{https://github.com/pgvector/pgvector}

\bibitem{assetmind}
\system, ``Source repository and frozen evaluation artifacts,'' 2026.
[Online]. Available:
\url{https://github.com/KumarNirupam1/AssestMind_AI}
\end{thebibliography}

\end{document}
